\chapter{Konzeption und Implementierung der Vergleichsszenarien}
\label{ch:implementierung}

Im Folgenden wird die praktische Durchführung der Case Study dokumentiert. Dazu werden die Bewertungskriterien erarbeitet und definiert und darauf aufbauend die Hypothesen aufgebaut des Evaluationsrahmen in Abschnitt \ref{sec:evaluationsrahmen}. Des Weiteren wird die Grundanwendung und der technologische Aufbau in Abschnitt \ref{sec:grundanwendung} dokumentiert. Im Abschnitt \ref{sec:frontend} wird der Aufbau des Frontend beschrieben. Folgend werden im Abschnitt \ref{sec:spring_backend} und \ref{sec:sap_backend} die Implementierung des jeweiligen Backends dokumentiert.
\newline
Basierend auf dem in Abschnitt \ref{sec:hypothesen} definierten Hypothesen werden zwei prototypische Referenzarchitekturen implementiert. Ziel ist es, eine empirische Datengrundlage zu schaffen, um insbesondere die Unterschiede in der Sicherheitsarchitektur (Token-Validierung, Autorisierungslogik) und der Developer Experience (Setup, Registrierungs-Flows) zu identifizieren.

Die Implementierung isoliert die sicherheitskritischen Komponenten vom Rest der Geschäftslogik, um Störfaktoren zu minimieren. Dabei stehen die Metriken Konfigurationsaufwand ("Time-to-First-Token"), Implementierungskomplexität ("Lines of Code") und der Grad der Automatisierung im Fokus.


\section{Ableitung des Evaluationsrahmens und der Bewertungskriterien}
\label{sec:evaluationsrahmen}

Um den Vergleich zwischen \ac{AWS} Cognito und \ac{SAP IAS} im Rahmen der Case Study nach \citep{Kitchenham1995} systematisch durchzuführen, werden objektive Bewertungskriterien definiert. Diese Kriterien werden unter anderem aus zwei etablierten Industriestandards abgeleitet, nämlich den \ac{OWASP} Cloud-Native Application Security Top 10 zur Identifikation von Risiken\footnote{\ref{fn:owaspcloudnative}} und der ISO/IEC 25010\footnote{\url{https://iso25000.com/en/iso-25000-standards/iso-25010}, aufgerufen am 23.11.2025} zur Klassifizierung von Softwarequalitätsmerkmalen.

\subsection{Vergleichsszenarios}

Dabei werden zwei spezifische Implementierungsszenarien untersucht:

\begin{enumerate}
    \item \textbf{SAP CAP mit SAP IAS:} Verwendung des SAP-Ökosystems mit deklarativer Sicherheitskonfiguration und automatischer Autorisierungsdurchsetzung auf Datenbankebene.

    \item \textbf{Spring Boot mit AWS Cognito:} Verwendung von Spring Boot und Spring Security mit imperativer Autorisierungskonfiguration.
\end{enumerate}

Dieser Vergleich ermöglicht es, nicht nur die funktionalen Unterschiede der \ac{IdP} zu messen, sondern auch die Auswirkungen des Framework-Ansatzes auf die praktische Sicherheit, Wartbarkeit und Developer Experience zu bewerten.


\subsection{Bedrohungsanalyse im Kontext von Cloud-Identitätsanbietern}
\label{sec:bedrohung}

Die Analyse der aktuellen Bedrohungslage zeigt, dass Sicherheitsrisiken in der Cloud häufig nicht auf Protokollschwächen basieren, sondern wie diese implementiert und konfiguriert werden \citep{Dietrich2018}. Die \ac{OWASP} zeigt zudem in der \textit{Cloud-Native Application Security Top 10} spezifische Risiken, die für auch das \ac{IAM} wichtig sind\footnote{\ref{fn:owaspcloudnative}}.

\begin{itemize}
    \item \textbf{Insecure Cloud Configuration (\ac{OWASP} CNAS-1):} Dies ist das aktuell größte Risiko in Cloud-Umgebungen. Im Kontext von \ac{IAM} bedeutet dies, dass die Sicherheit eines Systems oft davon abhängt, wie die Standardkonfigurationen gegen menschliche Fehler schützen. Ein zu hoher Grad an Flexibilität ohne Sicherheitsleitplanken kann die Wahrscheinlichkeit von Fehlkonfigurationen erhöhen \citep{Dietrich2018}. Daher ist auch relevant, wie das Framework Entwickler durch Standardeinstellungen schützt.

    \item \textbf{Improper Authentication \& Authorization (\ac{OWASP} CNAS-3):} Fehler in der Validierung von Identitäten entstehen häufig durch komplexe Implementierungen in der Anwendungslogik. Daher ist zu untersuchen, ob die Frameworks Mechanismen bereitstellen, die diese Komplexität durch Abstraktion verringern.
\end{itemize}

\subsection{Definition der Bewertungskriterien}
\label{sec:bewertungskriterien}

Um eine Vergleichbarkeit zwischen den beiden Ökosystemen herzustellen, werden messbare Bewertungskriterien definiert. Diese orientieren sich am Qualitätsmodell der ISO/IEC 25010 und berücksichtigen die Sicherheitsrisiken aus den \ac{OWASP} Cloud-Native Application Security Top 10.
\newline
Die Bewertungskriterien werden unter anderem aus den Anforderungen, die bei der praktischen Evaluation von \ac{IdP} relevant sind, abgeleitet. Aus der Bedrohungsanalyse in Abschnitt \ref{sec:bedrohung} wird deutlich, dass Sicherheitsrisiken in der Cloud häufig nicht durch Protokollmängel entstehen, sondern durch Implementierungsfehler und Fehlkonfigurationen \citep{Dietrich2018}. Daher muss eine \ac{IdP} nicht nur funktional korrekt sein, sondern auch durch das Design und die Frameworkintegration Entwickler vor häufigen Fehlern bewahren \citep{Green2016}. Daher wird das Kriterium \textbf{Sicherheit} mit den Unterkriterien Vertraulichkeit, Integrität und Zurechenbarkeit gewählt, da diese direkt die OWASP-Risiken \textit{Insecure Cloud Configuration} (CNAS-1) und \textit{Improper Authentication \& Authorization} (CNAS-3) adressieren.

Aus der technologischen Integration im Abschnitt \ref{sec:tech_frameworks} wird zudem angedeutet, dass unterschiedliche Framework-Paradigmen (deklarativ vs. imperativ) Auswirkungen auf die Implementierungserfahrung haben. Daher wird \textbf{Developer Experience} als Kriterium herangezogen, das sowohl die initiale Lernkurve (Erlernbarkeit) als auch operative Aspekte (Analysierbarkeit) berücksichtigt \citep{Myers2016}. Ein hoher Implementierungsaufwand erhöht nicht nur die Kosten, sondern auch das Fehlerrisiko, weshalb dieses Kriterium für die Gesamtbewertung entscheidend ist \citep{Green2016}.

Außerdem laufen Identitätsverwaltungssysteme über lange Zeiträume im Betrieb und zudem können die Anforderungen sich ändern, daher ist \textbf{Wartbarkeit} auch ein weiteres Kriterium zur Bewertung. Die Modularität der Autorisierungslogik und deren Wiederverwendbarkeit über mehrere Anwendungen hinweg beeinflussen direkt die Betriebskosten und die Stabilität des Systems bei Änderungen \citep{Myers2016}.

Die verbleibenden Kriterien (\textbf{Zuverlässigkeit}, \textbf{Übertragbarkeit} und \textbf{Interoperabilität}) adressieren die identifizierte Multi-Cloud-Realität \citep{Ali2025}. Die Zuverlässigkeit ist fundamental für Systeme, die als zentraler Authentifizierungspunkt fungieren. Übertragbarkeit und Interoperabilität sind zudem auch essentiell zur Vermeidung von \textit{Vendor Lock-in} und zur Integration in Unternehmenslandschaften \citep{Zhuravel2024}. Ein Vendor Lock-in bedeutet, dass z.B. ein Kunde an einen Dienstleister gebunden ist und nicht mehr einfach wechseln kann.


\subsubsection{Sicherheit}

Sicherheit wird als die Fähigkeit des Systems verstanden, Risiken durch Fehlkonfigurationen und manuelle Implementierungsfehler zu minimieren. Gemäß \ac{OWASP} ist dies oft wichtiger als die reine Protokollstärke, da die meisten Sicherheitsvorfälle auf menschliche Fehler bei der Konfiguration zurückgehen \citep{Dietrich2018}.

\begin{itemize}
    \item \textbf{Vertraulichkeit:} Sensible Daten wie Client Secrets, User Attributes und personenbezogene Informationen müssen vor unberechtigtem Zugriff geschützt sein. Dies gilt sowohl für Speicherung als auch für Übertragung von Authentifizierungsdaten.

    \item \textbf{Integrität:} Das System muss die Authentizität und Integrität von Authentifizierungsdaten gewährleisten. Wichtig dabei ist, ob das Framework die Validierung automatisch durchführt oder manuelle Implementierung erforderlich ist.

    \item \textbf{Zurechenbarkeit:} Administrative Änderungen an der \ac{IdP}-Konfiguration müssen aufgezeichnet werden. Dadurch wird eine Ursachenforschung ermöglicht im Schadensfall.

    \item \textbf{Recovery-Mechanismen:} Können Benutzer sich selbst wiederherstellen
    (Password-Vergessen-Flow) oder ist Admin-Intervention notwendig? Außerdem wie gut sind Fehlerszenarien (z.B. Passwort vergessen) implementiert? Erfordern diese manuelle
    Implementierung oder sind sie standardisiert?
\end{itemize}

\subsubsection{Developer Experience}

Dieses Kriterium beschreibt, wie einfach Entwickler die \ac{IdP} in ihre Anwendung integrieren können. Es kombiniert Erlernbarkeit und die Qualität von Fehlerdiagnose.

\begin{itemize}
    \item \textbf{Erlenbarkeit:} Der initiale Aufwand, um eine funktionsfähige und sichere Authentifizierung zu implementieren, gemessen als \textit{Time-to-First-Token}. Die Konfigurationskomplexität wirkt sich direkt auf die Dauer dieser Phase aus. Hohe Komplexität führt oft zu nicht optimalen Lösungen, die Sicherheitslücken ermöglichen.

    \item \textbf{Analysierbarkeit:} Die Qualität von Fehlermeldungen, Logging-Mechanismen und Dokumentation bei Authentifizierungsfehlern. Dies ist besonders in der Betriebsphase der Anwendung wichtig für Problemdiagnosen.
\end{itemize}

\subsubsection{Wartbarkeit}

Cloud-Sicherheitsrichtlinien und Geschäftsanforderungen können sich ändern. Daher wird bewertet, wie leicht sich Autorisierungsregeln anpassen lassen, ohne die Stabilität zu gefährden.

\begin{itemize}
    \item \textbf{Modularität:} Der Aufwand, um Änderungen an der Autorisierungslogik vorzunehmen (z.\,B. neue Rollen, neue Scopes, neue Richtlinien). Ein hohes Maß an Modularität reduziert das Fehlerrisiko bei Wartungsarbeiten.

    \item \textbf{Wiederverwendbarkeit:} Die Fähigkeit, denselben \ac{IdP} und Autorisierungskonfiguration für mehrere Anwendungen zu nutzen. Dies ist besonders in Unternehmensumgebungen mit mehreren Systemen relevant.
\end{itemize}

\subsubsection{Zuverlässigkeit}

Der \ac{IdP} stellt einen zentralen Einstiegspunkt dar. Seine Zuverlässigkeit ist entscheidend für die Gesamtsystemverfügbarkeit.

\begin{itemize}
    \item \textbf{Verfügbarkeit:} Der Authentifizierungsdienst muss erreichbar sein, wenn benötigt.

    \item \textbf{Fehlertoleranz:} Das Verhalten der Anwendung, wenn der \ac{IdP} temporär nicht erreichbar, durch z.B. Caching.
\end{itemize}

\subsubsection{Übertragbarkeit}

Multi-Cloud-Strategien erfordern die Vermeidung von Vendor Lock-in. Damit \ac{IdP} u. a. austauschbar bleiben oder man diese leicht anpassen kann.

\begin{itemize}
    \item \textbf{Austauschbarkeit:} Wie leicht kann der \ac{IdP} ausgetauscht werden? Entscheidend ist, ob die Implementierung standardisierte Protokolle (\ac{OIDC}/\ac{SAML}) nutzt oder eigene Implementierungen. Zudem bestimmt der Koppelungsgrad zwischen Applikationscode und Framework die Metrik.

    \item \textbf{Anpassungsfähigkeit:} Die Fähigkeit des Systems, sich an neue technische oder organisatorische Anforderungen anzupassen, ohne Neuimplementierung.
\end{itemize}

\subsubsection{Interoperabilität}

In Unternehmenslandschaften muss der \ac{IdP} mit bestehenden Infrastrukturen zusammenarbeiten.

\begin{itemize}
    \item \textbf{Föderationsunterstützung:} Die Fähigkeit, externe \ac{IdP} über standardisierte Protokolle, wie \ac{SAML} oder \ac{OIDC} anzubinden und als zentraler Identity Hub zu fungieren.

    \item \textbf{Standardkonformität:} Der Grad der Konformität mit offenen Standards gewährleistet langfristige Interoperabilität.
\end{itemize}

\subsection{Hypothesenbildung für die Case Study}
\label{sec:hypothesen}

Basierend auf den definierten Bewertungskriterien, der Bedrohungsanalyse und der Analyse der technologischen Integration werden folgende Hypothesen aufgestellt und in der Fallstudie getestet. Nach den Prinzipien der evaluativen Case Study nach \citep{Kitchenham1995} werden alle Hypothesen unter kontrollierten Bedingungen getestet, um andere Faktoren, wie z.B. Netzwerkleistung auszuschließen und eine faire Vergleichbarkeit zwischen \ac{AWS} Cognito und \ac{SAP CAP} mit \ac{SAP IAS} zu gewährleisten.
\newline
Die Hypothesen sind zudem in zwei Kategorien unterteilt, um einmal die Framework-Integration und auch die \ac{IdP}s selbst zu evaluieren. Die erste Kategorie ($H_1$ bis $H_5$) untersucht die Framework-Integration mit den jeweiligen \ac{IdP}s, also Spring Boot mit \ac{AWS} Cognito und \ac{SAP CAP} mit \ac{SAP IAS} die Authentifizierung und Autorisierung implementieren. Die zweite Kategorie ($H_6$ bis $H_9$) fokussiert sich auf den direkten Vergleich der \ac{IdP}s (\ac{AWS} Cognito vs. \ac{SAP IAS}), unabhängig von der Framework-Implementierung.

\subsubsection{Hypothesen zur Framework-Integration mit IdPs}

\begin{enumerate}
    \item[$H_1$] \textbf{Security:} Der deklarative Sicherheitsansatz von \ac{SAP CAP} reduziert die Anzahl notwendiger Autorisierungs-Code-Zeilen gegenüber dem imperativen Spring Security Ansatz mit \ac{AWS} Cognito. Dies senkt das Risiko für Implementierungsfehler wie \textit{Improper Authorization} (\ac{OWASP} CNAS-3).

    \item[$H_2$] \textbf{Maintainability:} Änderungen an der Autorisierungslogik (neue Rollen, Scopes, Richtlinien) erfordern im \ac{SAP CAP} mit \ac{SAP IAS} Ökosystem weniger Änderungen am Quellcode der Geschäftslogik als im Spring Boot mit \ac{AWS} Cognito Ökosystem, da die Autorisierungsregeln zentral im \ac{CDS}-Modell definiert werden.

    \item[$H_{3a}$] \textbf{Developer Experience:} \ac{SAP CAP} benötigt weniger Code und Konfiguration für Token-Validierung und \ac{JWT}-Handling als Spring Boot mit \ac{AWS} Cognito.

    \item[$H_{3b}$] \textbf{Developer Experience -- Unternehmensintegration:} Bei der Anbindung einer bestehenden Unternehmensidentität (mit externem \ac{IdP}) ist der Konfigurationsaufwand mit \ac{SAP CAP}/\ac{SAP IAS} in der Cloud geringer als mit Spring Boot mit \ac{AWS} Cognito in der Cloud.

    \item[$H_4$] \textbf{Portability:} Der Koppelungsgrad zwischen Applikationscode und Framework ist bei Spring Boot mit \ac{AWS} Cognito aufgrund der verteilten Spring Security Annotationen in den Service-Methoden höher als bei \ac{SAP CAP}/\ac{SAP IAS}, wo Autorisierungsregeln zentral im \ac{CDS}-Modell definiert sind.

    \item[$H_5$] \textbf{Security -- Token-Validierung:} Die Anzahl der Code-Zeilen für die \ac{JWT}-Token-Validierung ist bei \ac{SAP CAP} geringer als bei Spring Boot mit \ac{AWS} Cognito, da \ac{SAP CAP} die Validierung automatisch durch das Framework durchführt.
\end{enumerate}

\subsubsection{Hypothesen zum direkten IdP-Vergleich}

Die folgenden Hypothesen vergleichen \ac{AWS} Cognito und \ac{SAP IAS} als \ac{IdP}s direkt, ohne Berücksichtigung der Framework-Implementierung. Die Messungen erfolgen über die Admin-Konsolen und die integrierten Features der \ac{IdP}s.

\begin{enumerate}
    \item[$H_6$] \textbf{IdP -- Benutzererstellung:} Die manuelle Benutzererstellung über die Admin-Konsole erfordert bei \ac{SAP IAS} weniger Klicks/Schritte als bei \ac{AWS} Cognito.

    \item[$H_7$] \textbf{IdP -- Recovery-Mechanismen:} \ac{SAP IAS} bietet einen sichereren Passwort-Reset-Mechanismus als \ac{AWS} Cognito, da der Reset über einen zeitlich limitierten Link erfolgt statt über einen manuell einzugebenden Code.

    \item[$H_8$] \textbf{IdP -- Integrierte Sicherheitsfeatures:} \ac{SAP IAS} bietet mehr standardmäßig integrierte Sicherheitsfeatures (Rate-Limiting, Audit-Logging, automatische Token-Expiration, Brute-Force-Schutz) als \ac{AWS} Cognito.

    \item[$H_9$] \textbf{IdP -- Benutzerlöschung:} Die Löschung eines Benutzers über die Admin-Konsole erfordert bei \ac{AWS} Cognito weniger Klicks/Schritte als bei \ac{SAP IAS}.
\end{enumerate}

Die Hypothesen beschreiben zwei zentrale Ebenen des \ac{IAM}. Die \textbf{Framework-Integration} ($H_1$ bis $H_5$) testet die Kernunterscheidung zwischen deklarativen und imperativen Sicherheitsparadigmen sowie die Token-Validierung. Der \textbf{direkte IdP-Vergleich} ($H_6$ bis $H_9$) evaluiert die Entwickler-Usability und Sicherheitsfeatures der \ac{IdP}s selbst.

$H_1$ und $H_2$ testen die Kernunterscheidung zwischen deklarativen und imperativen Sicherheitsparadigmen, die das zentrale Differenzierungsmerkmal zwischen \ac{SAP CAP} und Spring Boot darstellt. $H_4$ adressiert das Risiko von Vendor Lock-in durch unterschiedliche Kopplungsmuster.

$H_5$ ist wichtig, da die automatische Token-Validierung durch \ac{SAP CAP} einen strukturellen Sicherheitsvorteil darstellt, der Implementierungsfehler (z.B. fehlende oder falsche Signaturverifizierung) verhindert. Dies ist insbesondere im Kontext des \ac{OWASP} Risikos ``\textit{Improper Authentication \& Authorization}'' zentral.

$H_{3a}$ und $H_{3b}$ messen die praktische Developer Experience in unterschiedlichen Szenarien. $H_{3a}$ misst die Code-Komplexität von Token-Mechanismen unter kontrollierten lokalen Bedingungen, während $H_{3b}$ die Unternehmensintegration in Cloud-Umgebungen validiert.

$H_6$ und $H_9$ vergleichen die Benutzerfreundlichkeit der Admin-Konsolen beider \ac{IdP}s durch objektive Klick-Messungen. $H_7$ evaluiert die Sicherheit der Recovery-Mechanismen (Link-basiert vs. Code-basiert). $H_8$ quantifiziert die standardmäßig integrierten Sicherheitsfeatures beider \ac{IdP}.


\subsubsection{Messbarkeiten der Hypothesen}

Die folgenden Hypothesen werden durch Metriken gemessen,
um ihre Überprüfbarkeit sicherzustellen. Jede Hypothese wird anhand von messbaren Kriterien evaluiert, die während der Implementierungsphase erfasst werden.
Die Hypothesen zeigen die Unterschiede zwischen den
Frameworks und \ac{IdP}s und sind nicht verschiedenen Entwickler-Kompetenzen abhängig. Alle Messungen basieren daher auf objektiven Metriken wie Code-Zeilen,
Anzahl von Dateien, zyklomatische Komplexität und Klick-Zählungen, die unabhängig von der Entwickler-Kompetenz sind.

\paragraph{Messbarkeiten für Framework-Integration ($H_1$ bis $H_5$)}

\textbf{H1 – Sicherheit der Autorisierungslogik:} Die Anzahl der Dateien mit Autorisierungslogik wird gezählt, um die Streuung der Sicherheitsregeln zu ermitteln.
Zusätzlich wird die Anzahl geschriebener Codezeilen
für Autorisierung gemessen. Für \ac{SAP CAP} werden alle \texttt{@restrict}-Annotationen
im \ac{CDS}-Modell gezählt. Für Spring Boot werden
alle \texttt{@PreAuthorize}-Annotationen über Controller und Services gezählt.

\textbf{H2 – Code-Änderungen und Wartbarkeit:} Durch das Hinzufügen einer neuen
Rolle wird gemessen, wie viele Dateien und Methoden angepasst werden müssen.

\textbf{H3a – Token-Validierung und JWT-Handling:} Die Anzahl der Code- und Konfigurationszeilen
für Token-Validierung und \ac{JWT}-Handling wird gemessen. Zusätzlich wird die zyklomatische Komplexität der Token-Validierungslogik mittels statischer Code-Analyse erfasst. Die zyklomatische Komplexität misst die Anzahl unabhängiger Pfade durch den Code und zeigt damit die Fehleranfälligkeit. Für \ac{SAP CAP} werden die Konfigurationszeilen in \texttt{application.yaml} gezählt. Für Spring Boot wird die \texttt{application.yaml}, SecurityConfig und JwtAuthenticationFilter analysiert.

\textbf{H3b – Unternehmensintegration:} Der
Konfigurationsaufwand bei der Anbindung eines externen \ac{IdP} wird erfasst
durch die Messung von erforderlichen Konfigurationsschritten und Codezeilen.

\textbf{H4 – Vendor-Abhängigkeit:} Die Anzahl von Dateien mit Framework- und
\ac{IdP}-spezifischen Importen wird gezählt. Dies misst, wie sehr der jeweilige \ac{IdP} gekoppelt ist. Für \ac{SAP CAP} werden Dateien mit @sap
spezifischen Imports gezählt. Für Spring Boot werden Dateien mit \texttt{org.springframework.security} Imports gezählt.

\textbf{H5 – Token-Validierung:} Es wird gemessen, ob die Signaturverifizierung
automatisch oder manuell erfolgt. Die zyklomatische Komplexität der Token-Validierungslogik wird mittels statischer Code-Analyse gemessen. Für \ac{SAP CAP} beträgt diese 0, da das Framework die Validierung automatisch durchführt.
Für Spring Boot wird die Komplexität des \texttt{JwtAuthenticationFilter} und \texttt{CognitoService.validateToken()} gemessen.

\paragraph{Messbarkeiten für direkten IdP-Vergleich ($H_6$ bis $H_9$)}

\textbf{H6 – Benutzererstellung:} Die manuelle Benutzererstellung wird in beiden Admin-Konsolen durch Anzahl der Klicks/Schritte gemessen:
\begin{itemize}
    \item Für \ac{AWS} Cognito: Anzahl der Klicks in der AWS Console für die Erstellung eines neuen Benutzers (User Pool $\rightarrow$ Users $\rightarrow$ Create User $\rightarrow$ Formular ausfüllen).
    \item Für \ac{SAP IAS}: Anzahl der Klicks in der SAP IAS Admin-Konsole für die Erstellung eines neuen Benutzers (Users \& Authorizations $\rightarrow$ User Management $\rightarrow$ Add $\rightarrow$ Formular ausfüllen).
\end{itemize}

\textbf{H7 – Recovery-Mechanismen:}
Der Passwort-Vergessen-Flow wird analysiert und anhand folgender Aspekte gemessen:
\begin{itemize}
    \item \textbf{Sicherheitsaspekt:} Gibt es einen Link-basierten oder Code-basierten Flow? Link-basierte Resets gelten als sicherer, da der Code nicht manuell übertragen werden muss.
    \item \textbf{Gültigkeitsdauer:} Zeitliche Limitierung des Reset-Tokens/-Codes in beiden \ac{IdP}s.
    \item \textbf{Konfigurationsaufwand:} Anzahl der Einstellungen im jeweiligen \ac{IdP}, um den Recovery-Flow zu aktivieren.
\end{itemize}

\textbf{H8 – Integrierte Sicherheitsfeatures:}
Die Anzahl der integrierten Sicherheitsfeatures wird für beide \ac{IdP} gezählt:
\begin{itemize}
    \item \textbf{Zu prüfende Features:} Rate-Limiting, Brute-Force-Schutz, Audit-Logging, automatische Token-Expiration, E-Mail-Verifizierung, Passwort-Policies, Account-Lockout, MFA-Unterstützung.
    \item Für jedes Feature wird dokumentiert: (a) standardmäßig aktiviert, (b) konfigurierbar, (c) nicht verfügbar.
\end{itemize}

\textbf{H9 – Benutzerlöschung:}
Die Löschung eines Benutzers wird in beiden Admin-Konsolen durch folgende Metriken gemessen:
\begin{itemize}
    \item Für \ac{AWS} Cognito: Anzahl der Klicks in der AWS Console für die Löschung eines Benutzers (User Pool $\rightarrow$ Users $\rightarrow$ Benutzer auswählen $\rightarrow$ Delete).
    \item Für \ac{SAP IAS}: Anzahl der Klicks in der SAP IAS Admin-Konsole für die Löschung eines Benutzers (Users \& Authorizations $\rightarrow$ User Management $\rightarrow$ Benutzer auswählen $\rightarrow$ Delete).
    \item Zudem wird die Anzahl der Bestätigungsschritte für die Löschung gezählt.
\end{itemize}

\subsubsection{Testumgebungen und Messkontext}

Die Hypothesen zur Framework-Integration ($H_1$ bis $H_5$) messen Framework Eigenschaften wie Code-Struktur, Autorisierungslogik und Vendor-Kopplung. Diese werden durch statische Code-Analyse erfasst und sind unabhängig von der Laufzeitumgebung.

Für die IdP-spezifischen Hypothesen ($H_6$ bis $H_9$) werden die beiden Ökosysteme in unterschiedlichen Umgebungen betrieben:

\begin{itemize}
    \item \textbf{Spring Boot mit \ac{AWS} Cognito:} Die Anwendung läuft lokal und verbindet sich direkt mit dem \ac{AWS} Cognito Cloud-Service. Das \ac{AWS} \ac{SDK} ermöglicht eine vollständige \ac{IdP}-Integration ohne Cloud-Deployment des Backends.

    \item \textbf{\ac{SAP CAP} mit \ac{SAP IAS}:} Da \ac{SAP IAS} ausschließlich als Cloud-Service verfügbar ist, wird das \ac{SAP CAP} Backend auf die \ac{SAP BTP} deployt. Die Anbindung an \ac{SAP IAS} erfolgt über automatisches Service-Binding. Lokal kann die Anwendung ohne \ac{IdP}-Anbindung gestartet werden, wobei die Framework-Architektur und die deklarativen Sicherheitsregeln weiterhin messbar sind.
\end{itemize}

Die Metriken werden während der Implementierung erfasst und in
einer Tabelle dokumentiert, um einen direkten Vergleich zwischen
\ac{AWS} Cognito mit Spring Boot und \ac{SAP IAS} mit \ac{SAP CAP}
zu ermöglichen.

\section{Definition der Grundanwendung}
\label{sec:grundanwendung}
Um die Vergleichbarkeit der beiden \ac{IdP}-Lösungen zu gewährleisten, wird eine Web-Anwendung in beiden Ökosystemen genutzt und entwickelt. Die Architektur folgt einer strikten Trennung von Frontend und Backend.

\subsection{Funktionale Anforderungen und Vergleichspunkte}
Die Anwendung setzt dabei Anforderungen um, die für die Prüfung der Hypothesen benötigt werden. Die Anforderungen gliedern sich in zwei Kategorien:

\paragraph{Backend-Implementierung (Framework-Integration $H_1$ bis $H_5$):}
\begin{itemize}
    \item \textbf{Benutzer-Verwaltung ($H_{3a}$):} Ein Benutzer soll sich anmelden können.
    \item \textbf{Authentifizierung und Token-Validierung ($H_5$):} Der Login-Prozess muss ein valides \ac{JWT} liefern.
    \item \textbf{Geschützter API-Zugriff ($H_1$, $H_2$):} Ein Backend-Endpunkt (z. B. \texttt{/api/orders}) wird rollenbasiert abgesichert.
\end{itemize}

\paragraph{IdP-Vergleich über Admin-Konsolen ($H_6$ bis $H_9$):}
Die folgenden Funktionen werden nicht im Frontend implementiert, sondern direkt über die Admin-Konsolen der jeweiligen \ac{IdP}s gemessen:
\begin{itemize}
    \item \textbf{Benutzerregistrierung ($H_6$):} Ein neuer Benutzer wird über die Admin-Konsole des jeweiligen \ac{IdP} erstellt. Die Anzahl der Klicks wird verglichen.
    \item \textbf{Recovery-Mechanismus ($H_7$):} Der Passwort-Vergessen-Flow wird verglichen (Link-basiert vs. Code-basiert). Für das Spring Boot Backend wird der Password-Vergessen-Flow zusätzlich implementiert.
    \item \textbf{Integrierte Sicherheitsfeatures ($H_8$):} Die Anzahl der integrierten Sicherheitsfeatures wird für beide \ac{IdP} verglichen und dokumentiert.
    \item \textbf{Benutzerlöschung ($H_9$):} Ein Benutzer wird über die Admin-Konsole des jeweiligen \ac{IdP} gelöscht. Die Anzahl der Klicks wird verglichen.
\end{itemize}

\subsection{Technologischer Aufbau}
Das Frontend dient als Kontrollvariable, um die Backend-Unterschiede messen zu können. Die Frontend-Komplexität ist daher gering gehalten, damit sich die Ergebnisse der Case Study auf die Backend-seitige Implementierung von Authentifizierung und Autorisierung fokussieren.
\textbf{Angular}\footnote{\url{https://angular.io/}, aufgerufen am 02.12.2025}
ist ein von Google entwickeltes Framework zur Erstellung von Single-Page-Webanwendungen
und wird in dieser Fallstudie als Plattform für einen standardisierten \ac{OIDC}-Client
verwendet. Die Wahl von Angular als konstante Frontend-Komponente dient dazu,
Frontend-spezifische Unterschiede als Einflussfaktor auszuschließen und die
Messbarkeit der Backend Unterschiede zu zeigen. Um Frontend-spezifische Unterschiede auszuschließen, wird das Frontend nach dem Backend-for-Frontend Pattern konzipiert.

\begin{figure}[h]
    \centering
    \includegraphics[width=1.\textwidth]{figures/bffarchitektur.drawio.pdf}
    \caption{Entwickelte Architektur anhand des Backend-for-Frontend Pattern}
    \label{fig:bffarch}
\end{figure}

Das Diagramm \ref{fig:bffarch} zeigt die Architektur, die in der Entwicklung umgesetzt wird. Dabei benutzt der Nutzer nur das Frontend, welches als HTTP-Client fungiert. Die Anfragen werden dann an das jeweilige Backend weitergeleitet. Die Backends rufen jeweils den gewählten \ac{IdP} auf. Somit passiert die gesamte Logik im Backend.
Dies ermöglicht es, die Unterschiede bei der Implementierung zwischen
\ac{AWS} Cognito und \ac{SAP IAS} isoliert zu messen.
Die Verwendung von Angular als konstante Frontend-Komponente dient der
systematischen Kontrolle von Variablen, wie Authentifizierungs-Flow, Token-Handling,
Benutzerführung und API-Kommunikation sind identisch in beiden Szenarien.
Dies stellt sicher, dass gemessene Unterschiede ausschließlich auf Backend-seitige
Implementierungen von \ac{AWS} Cognito und \ac{SAP IAS} zurückgehen und nicht
durch Frontend-Designentscheidungen beeinflusst werden.

\begin{itemize}
    \item \textbf{Frontend (Konstante):} Eine \textbf{Angular}-Anwendung (aktuelle Version), die als generischer \ac{OIDC}-Client agiert. Sie übernimmt die Benutzerführung und leitet Tokens an das Backend weiter.
    \item \textbf{Backend (Variable):} Es werden zwei separate Backends implementiert, ein Spring Boot Backend mit (\ac{AWS} SDK) und ein Backend basierend auf \ac{SAP CAP} Framework.
\end{itemize}

\section{Angular Frontend}
\label{sec:frontend}

Das Frontend basiert auf dem Angular-Framework als konstante Komponente zur Umsetzung eines standardisierten \ac{OIDC}-Clients. Dadurch wird sichergestellt, dass Frontend-spezifische Unterschiede als Einflussfaktor ausgeschlossen werden und die Messbarkeit der Backend-Unterschiede gegeben ist. Die Architektur folgt dabei dem Backend-for-Frontend Pattern, wobei der Nutzer nur das Frontend als HTTP-Client nutzt und Anfragen an die jeweiligen Backends weitergeleitet werden.


\begin{figure}[h]
    \centering
    \includegraphics[width=1.\textwidth]{figures/frontend-login.png}
    \caption{Screenshot der Login-Seite mit Backend Auswahl}
    \label{fig:frontend-login}
\end{figure}

Abbildung \ref{fig:frontend-login} zeigt die Login Oberfläche als Beispiel. Dabei gibt es oberhalb eine Auswahl-Box, die die Wahl zwischen den zwei Backends ermöglicht und außerdem wird angezeigt, welches Backend aktuell genutzt wird. Darunter sind zwei Input-Felder für die Eingabe des Benutzernamens und des Passworts angeordnet. Unterhalb befinden sich zwei Buttons, einer führt den Authentifizerungsflow durch und der zweite Button wechselt zur Registrierungsoberfläche.

\subsubsection{Implementierte Funktionalitäten}

Neben der Login-Seite wurden folgende Komponente und Services für das Frontend entwickelt:

\begin{itemize}
    \item \textbf{Registrierungs-Komponente:} Die Registrierungsseite ist nur für das Spring Boot Backend mit \ac{AWS} Cognito verfügbar. Bei Auswahl des \ac{SAP CAP} Backends wird stattdessen ein Hinweis angezeigt, dass die Benutzerregistrierung über die \ac{SAP IAS} Admin-Konsole erfolgt. Dies zeigt den Unterschied in der Architektur, denn \ac{AWS} Cognito ermöglicht auch programmatische Registrierung, während \ac{SAP IAS} eine administrative Benutzererstellung erfordert.
    \item \textbf{Dashboard-Komponente:} Nach erfolgreicher Authentifizierung wird der Nutzer automatisch auf ein Dashboard weitergeleitet, das einen Überblick über seine Informationen und verfügbaren Funktionen bietet.
    \item \textbf{Geschützte Ressourcenseite:} Nach erfolgreichem Login können Nutzer auf geschützte Ressourcen zugreifen. Im Frontend wird dies exemplarisch durch eine Order-Seite gezeigt, die nur authentifizierten Nutzern mit entsprechenden Berechtigungen/festgelegte Rollen aufrufbar ist.
    \item \textbf{Callback-Komponente:} Eine Callback-Komponente verarbeitet die Rückkehr vom Authorization Server nach erfolgreichem Login und leitet den Nutzer automatisch auf das Dashboard um.
    \item \textbf{Passwort-Vergessen-Komponente:} Eine Seite, wo der Nutzer seine E-Mail eintragen kann und anschließend das Ergebnis von der zugesendeten E-Mail, um den Flow abzuschließen.
    \item \textbf{Service-Methoden für Backend-Kommunikation:} Zudem wurden Service-Methoden definiert, um die HTTP-Kommunikation mit den Backend zu ermöglichen.
\end{itemize}

\subsubsection{Authentifizierungsflow}

Der Authentifizierungsflow erfolgt nach dem Standard-\ac{OIDC} Authorization Code Flow. Dabei wählt der Nutzer sein Backend, füllt seine Anmeldedaten ein und klickt auf Login. Das Frontend leitet die Anfrage an das Backend weiter, das wiederum den konfigurierten \ac{IdP} (\ac{AWS} Cognito oder \ac{SAP IAS}) aufruft. Der \ac{IdP} authentifiziert den Nutzer und gibt einen Authorization Code zurück. Das Backend tauscht diesen Code gegen ein \ac{JWT} aus und gibt dieses an das Frontend zurück. Das Frontend speichert das Token lokal und nutzt es für alle weiteren API-Anfragen. Dadurch wird gewährleistet, dass die gesamte Sicherheitslogik im Backend verbleibt und Frontend-Designentscheidungen die Backend-Unterschiede nicht beeinflussen.

Die Service-Methoden abstrahieren dabei die jeweiligen API-Endpunkte und verwalten zentral das Token-Handling, wodurch sich die Anwendungslogik nicht um Backend-spezifische Unterschiede kümmern muss. Diese Entkopplung ermöglicht es, dass beide Backend-Implementierungen von derselben Frontend-Anwendung genutzt werden können, ohne dass Codenanpassungen im Frontend selbst notwendig sind.


\section{Spring Boot mit AWS SDK Backend und AWS Cognito}
\label{sec:spring_backend}

Im Folgenden wird die Integration von \ac{AWS} Cognito unter Verwendung des Spring Boot Frameworks dokumentiert. Dabei liegt der Fokus auf der Architektur, der Konfiguration sowie den zentralen Komponenten für Authentifizierung und Autorisierung. Die Implementierung nutzt das \ac{AWS} \ac{SDK} für Java in Kombination mit Spring Security.

\subsection{Architektur und Projektstruktur}

Das Spring Boot Backend folgt einer klassischen Schichtenarchitektur mit Trennung von Controller-, Service- und Konfigurationsschichten. Die Projektstruktur ist dabei wie folgt implementiert:

\begin{itemize}
    \item \textbf{controller}: Enthält die REST-Controller für Authentifizierung (\texttt{AuthController}) und Geschäftslogik (\texttt{OrderController})
    \item \textbf{service}: Beinhaltet den \texttt{CognitoService} für die Kommunikation mit \ac{AWS} Cognito sowie den \texttt{OrderService} für die Geschäftslogik
    \item \textbf{security}: Enthält den \texttt{JwtAuthenticationFilter} für die manuelle Token-Validierung
    \item \textbf{config}: Umfasst die \texttt{SecurityConfig} für Spring Security und die \texttt{AwsCognitoProperties} für die Konfiguration
    \item \textbf{dto}, \textbf{model}, \textbf{exception}: Unterstützende Klassen für Datenübertragung, Entitäten und Fehlerbehandlung
\end{itemize}

Diese Struktur entspricht dem imperativen Programmiermodell, bei dem Sicherheitsregeln und Geschäftslogik sind im Code implementiert.

\subsection{AWS Cognito User Pool Einrichtung}

Bevor die Spring Boot Anwendung konfiguriert werden kann, wird zunächst ein User Pool in \ac{AWS} erstellt und konfiguriert. Der User Pool dient als \ac{IdP} für die Anwendung.

Die Einrichtung umfasst mehrere Konfigurationsschritte. Zunächst wird ein neuer User Pool mit einem eindeutigen Namen erstellt. Dabei werden die Passwort-Richtlinien festgelegt (Mindestlänge, Zahlen, usw.). Für diese Implementierung wird eine Mindestlänge von 8 Zeichen mit allen Zeichenklassen gewählt, da die Anwendung nur lokal läuft, im produktiv Einsatz wäre die gewählte Passwort-Richtlinie nicht sicher.

Anschließend wird ein App Client innerhalb des User Pools erstellt. Abbildung \ref{fig:cognito_appclient} zeigt die Konfiguration des App Clients in \ac{AWS}. Der App Client definiert, welche Authentifizierungs-Flows die Anwendung nutzen darf. Für diese Implementierung wurden u. a. folgende Einstellungen vorgenommen:
\begin{itemize}
    \item \textbf{Client ID und Client Secret:} Generierte Credentials für die sichere Kommunikation zwischen Backend und Cognito
    \item \textbf{Auth Flow:} \texttt{ALLOW\_USER\_PASSWORD\_AUTH} für den direkten Passwort-basierten Login
    \item \textbf{Token-Gültigkeit:} Access Token 1 Stunde, Refresh Token 5 Tage
    \item \textbf{E-Mail-Verifizierung:} Aktiviert mit automatischem Versand eines 6-stelligen Codes
\end{itemize}

\begin{figure}[h]
    \centering
    \includegraphics[width=1.1\textwidth]{figures/userpool.png}
    \caption{Konfiguration des App Clients in AWS Cognito}
    \label{fig:cognito_appclient}
\end{figure}

Die im App Client generierten Werte (User Pool ID, Client ID, Client Secret) werden anschließend in der \texttt{application.properties} des Spring Boot Backends hinterlegt.

\subsection{Konfiguration und Einrichtung}

Die Konfiguration erfolgt primär über die \texttt{application.properties}-Datei, in der alle \ac{AWS} Cognito-spezifischen Parameter definiert werden. Dies umfasst die User Pool ID, Client ID, Client Secret sowie die \ac{JWKS}-\ac{URI} für die Token-Validierung. Abbildung \ref{fig:spring_properties} zeigt die wesentlichen Konfigurationsparameter.

\begin{figure}[h]
    \begin{verbatim}
# AWS Cognito Configuration
aws.region=eu-central-1
aws.cognito.userPoolId=eu-central-1_QuhDIcGXQ
aws.cognito.clientId=2p1c8sb12boir1dkid42sfc5bh
aws.cognito.clientSecret=***
aws.cognito.jwksUri=https://cognito-idp.${aws.region}.amazonaws.com/
    ${aws.cognito.userPoolId}/.well-known/jwks.json
aws.cognito.issuerUri=https://cognito-idp.${aws.region}.amazonaws.com/
    ${aws.cognito.userPoolId}

spring.security.oauth2.resourceserver.jwt.issuer-uri=${aws.cognito.issuerUri}
spring.security.oauth2.resourceserver.jwt.jwk-set-uri=${aws.cognito.jwksUri}
    \end{verbatim}
    \caption{Konfiguration der AWS Cognito Parameter in \texttt{application.properties}}
    \label{fig:spring_properties}
\end{figure}

Wie in Abbildung \ref{fig:spring_properties} gezeigt, müssen die \ac{AWS}-spezifischen Parameter als auch die Spring Security OAuth2-Konfiguration definiert werden. Die \texttt{issuer-uri} und \texttt{jwk-set-uri} werden dabei aus den \ac{AWS} Cognito-Parametern abgeleitet, um die \ac{JWT}-Validierung zu ermöglichen.

Zusätzlich werden die Konfigurationsparameter über eine Properties-Klasse (\texttt{AwsCognitoProperties}) in die Anwendung geladen, was die Wiederverwendung in verschiedenen Services ermöglicht.

\subsection{Security-Konfiguration}

Die Sicherheitskonfiguration erfolgt über die Klasse \texttt{SecurityConfig}, die Spring Security für die \textit{stateless} \ac{JWT}-basierte Authentifizierung konfiguriert. Stateless Applikationen speichern keine Nutzerdaten serverseitig, stattdessen werden alle benötigten Informationen bei jeder Anfrage erneut mitgesendet\footnote{\url{https://www.redhat.com/en/topics/cloud-native-apps/stateful-vs-stateless}, aufgerufen am 25.12.2025}. Die Konfiguration definiert öffentliche und geschützte Endpunkte sowie die Integration des eigenen \texttt{JwtAuthenticationFilter}.

Die \texttt{SecurityConfig} aktiviert die methodenbasierte Sicherheit durch \texttt{@EnableMethodSecurity(prePostEnabled = true)}, wodurch Annotationen wie \texttt{@PreAuthorize} auf Controller-Methoden verwendet werden können. Zudem wird die Session-Verwaltung auf \texttt{STATELESS} gesetzt, da keine serverseitigen Sessions verwendet werden. Der eigene \texttt{JwtAuthenticationFilter} wird vor dem Standard-Authentifizierungsfilter von Spring Security eingefügt, um eingehende \ac{JWT}-Tokens zu validieren.

\subsection{CognitoService: Authentifizierung und Token-Validierung}

Der \texttt{CognitoService} ist die zentrale Komponente für die Kommunikation mit \ac{AWS} Cognito. Er implementiert alle Authentifizierungs- und Registrierungsflows und nutzt das \ac{AWS} \ac{SDK} für Java. Abbildung \ref{fig:cognito_service} zeigt exemplarisch die manuelle \ac{JWT}-Validierung.

\begin{figure}[h]
    \begin{verbatim}
public JWTClaimsSet validateToken(String token) {
        ...
        String issuer = claims.getIssuer();
        String expectedIssuer = cognitoProperties.getCognito().getIssuerUri();
        if (!expectedIssuer.equals(issuer)) {
            throw new InvalidTokenException("Invalid token issuer");
        }

        Instant expirationTime = claims.getExpirationTime().toInstant();
        if (expirationTime.isBefore(Instant.now())) {
            throw new TokenExpiredException("Token has expired");
        }
        ...
}
    \end{verbatim}
    \caption{Ausschnitt aus der Methode \texttt{validateToken(...)} für manuelle JWT-Validierung im CognitoService}
    \label{fig:cognito_service}
\end{figure}

Wie in Abbildung \ref{fig:cognito_service} dargestellt, erfordert die Token-Validierung eine manuelle Implementierung. Die Validierung umfasst dabei das Laden des öffentlichen Schlüssels vom \ac{JWKS}-Endpunkt, die Signaturprüfung, sowie die Validierung von Issuer, Expiration und Audience. Falls der Issuer nicht übereinstimmt oder der Token abgelaufen ist wird eine entsprechende Exception geworfen.

Der \texttt{CognitoService} implementiert zusätzlich folgende Funktionen:
\begin{itemize}
    \item \textbf{Benutzerregistrierung:} Registrierung neuer Benutzer über das \ac{AWS} \ac{SDK} mit Attributzuweisung (E-Mail, Vorname, Nachname)
    \item \textbf{E-Mail-Verifizierung:} Bestätigung der Registrierung mit einem 6-stelligen Verifizierungscode
    \item \textbf{Login:} Authentifizierung über den \texttt{USER\_PASSWORD\_AUTH}-Flow mit Rückgabe von Access Token, ID Token und Refresh Token
    \item \textbf{Passwort zurücksetzen:} Implementierung des Passwort-Vergessen-Flows mit Code-Verifizierung ($H_7$)
    \item \textbf{Secret Hash Berechnung:} Berechnung des HMAC-SHA256-Hashes für Clients mit Client Secret
\end{itemize}

\subsection{AuthController: REST-Endpunkte für Authentifizierung}

Der \texttt{AuthController} stellt die REST-\ac{API} für alle Authentifizierungsoperationen bereit. Zudem leitet er Anfragen an den \texttt{CognitoService} weiter und gibt das Ergebnis zurück.

Die implementierten Endpunkte umfassen:
\begin{itemize}
    \item \texttt{POST /api/auth/register}: Registrierung neuer Benutzer
    \item \texttt{POST /api/auth/login}: Authentifizierung und Token-Ausstellung
    \item \texttt{POST /api/auth/verify-email}: E-Mail-Verifizierung nach Registrierung
    \item \texttt{POST /api/auth/resend-verification-code}: Erneutes Senden des Verifizierungscodes
    \item \texttt{POST /api/auth/forgot-password}: Anforderung eines Passwort-Reset-Codes
    \item \texttt{POST /api/auth/confirm-password-reset}: Bestätigung des Passwort-Resets
\end{itemize}

Alle Endpunkte unter \texttt{/api/auth/**} sind öffentlich zugänglich, da sie vor der Authentifizierung aufgerufen werden müssen.

Der Passwort-Vergessen-Flow erfordert zudem eine zweistufige Implementierung. Im ersten Schritt ruft der Benutzer den \texttt{forgot-password}-Endpunkt mit seiner E-Mail-Adresse auf. Der \texttt{CognitoService} sendet daraufhin über die \ac{AWS} Cognito \ac{API} einen 6-stelligen Verifizierungscode an die registrierte E-Mail. Im zweiten Schritt übermittelt der Benutzer diesen Code zusammen mit dem neuen Passwort an den \texttt{confirm-password-reset}-Endpunkt.

\subsection{OrderController: Rollenbasierte Autorisierung}

Der \texttt{OrderController} setzt die imperative Autorisierung mit Spring Security Annotationen um. Abbildung \ref{fig:order_controller} zeigt die Verwendung der \texttt{@PreAuthorize}-Annotation.

\begin{figure}[h]
    \begin{verbatim}
    ...
    @GetMapping
    @PreAuthorize("isAuthenticated()")
    public ResponseEntity<List<Order>> getAllOrders() {
        List<Order> orders = orderService.getAllOrders();
        return ResponseEntity.ok(orders);
    }
    ...
    @PostMapping
    @PreAuthorize("hasRole('ADMIN')")
    public ResponseEntity<Order> createOrder(@RequestBody OrderDto orderDto) {
        Authentication auth = SecurityContextHolder.getContext().getAuthentication();
        Order order = orderService.createOrder(orderDto, auth.getName());
        return ResponseEntity.status(HttpStatus.CREATED).body(order);
    }
}
    \end{verbatim}
    \caption{Ausschnitt aus der \texttt{OrderController.java}-Datei mit Autorisierung mit @PreAuthorize}
    \label{fig:order_controller}
\end{figure}

Wie in Abbildung \ref{fig:order_controller} gezeigt, wird die Autorisierung durch \texttt{@PreAuthorize}-Annotationen auf jeder Methode definiert. Die einzelnen Methoden, können dadurch jeweils nur mit der geeigneten Rolle aufgerufen werden, wie beispielsweise \texttt{createOrder} mit der Rolle ADMIN. Dies entspricht dem imperativen Ansatz von Spring Security, bei dem Autorisierungsregeln über den gesamten Controller verteilt sind. Die Autorisierungslogik unterscheidet zwischen:
\begin{itemize}
    \item \texttt{isAuthenticated()}: Jeder authentifizierte Benutzer hat Zugriff
    \item \texttt{hasRole('MANAGER')}: Nur Benutzer mit der MANAGER-Rolle oder höher haben Zugriff
    \item \texttt{hasRole('ADMIN')}: Nur Benutzer mit der ADMIN-Rolle haben Zugriff
\end{itemize}

Die Rollen werden aus dem \ac{JWT}-Token extrahiert, durch das Attribut \texttt{cognito:groups}-Claim, und im \texttt{JwtAuthenticationFilter} in Spring Security Authorities umgewandelt. Die Rollen sind zudem hierarchisch aufgebaut, dabei ist die Rolle \texttt{Admin} die stärkste Rolle und \texttt{User} die schwächste.

\subsection{JwtAuthenticationFilter: Token-Extraktion und Authentifizierung}

Der \texttt{JwtAuthenticationFilter} ist ein Spring Security Filter, der eingehende HTTP-Anfragen auf \ac{JWT}-Tokens prüft und die Authentifizierung im Security Context setzt. Er extrahiert das Token aus dem \texttt{Authorization}-Header, validiert es über den \texttt{CognitoService} und wandelt die Cognito-Gruppen in Spring Security Authorities um.

Die Extraktion der Authorities erfolgt dabei aus dem \texttt{cognito:groups}-Claim des Tokens. Jede Gruppe wird mit dem Präfix \texttt{ROLE\_} versehen, um mit Spring Security Annotationen wie \texttt{hasRole('MANAGER')} kompatibel zu sein.

\subsection{Zusammenfassung der Implementierung}

Die Implementierung des Spring Boot Backends mit \ac{AWS} Cognito stellt den imperativen Sicherheitsansatz des Frameworks dar. Dabei gibt es folgende zentrale Aspekte:

\begin{itemize}
    \item \textbf{Konfiguration (\texttt{application.properties}):} Die Definition aller \ac{AWS} Cognito-spezifischen Parameter wie User Pool ID, Client ID, Client Secret sowie die \ac{JWKS}-\ac{URI} und Issuer-\ac{URI} für die Token-Validierung (vgl. Abbildung \ref{fig:spring_properties}).

    \item \textbf{Security-Konfiguration (\texttt{SecurityConfig.java}):} Die Spring Security Konfiguration für \textit{stateless} \ac{JWT}-basierte Authentifizierung, Definition öffentlicher und geschützter Endpunkte sowie die Integration des \texttt{JwtAuthenticationFilter}.

    \item \textbf{CognitoService (\texttt{CognitoService.java}):} Die zentrale Komponente für die Kommunikation mit \ac{AWS} Cognito über das \ac{AWS} \ac{SDK}. Implementiert die manuelle \ac{JWT}-Validierung, Benutzerregistrierung, Login und Passwort-Reset (vgl. Abbildung \ref{fig:cognito_service}).

    \item \textbf{JwtAuthenticationFilter (\texttt{JwtAuthenticationFilter.java}):} Der Spring Security Filter für die Extraktion und Validierung von \ac{JWT}-Tokens aus dem \texttt{Authorization}-Header sowie die Umwandlung von Cognito-Gruppen in Spring Security Authorities.

    \item \textbf{AuthController (\texttt{AuthController.java}):} Die REST-\ac{API} für alle Authentifizierungsoperationen (Register, Login, Verify-Email, Forgot-Password, Confirm-Password-Reset).

    \item \textbf{OrderController (\texttt{OrderController.java}):} Die Implementierung der rollenbasierten Autorisierung mit \texttt{@PreAuthorize}-Annotationen auf jeder Methode (vgl. Abbildung \ref{fig:order_controller}).
\end{itemize}

\paragraph{Relevanz für die Hypothesen}

\begin{itemize}
    \item \textbf{Manuelle Token-Validierung ($H_{3a}$, $H_5$):} Die \ac{JWKS}-\ac{URI} und Issuer-\ac{URI} müssen explizit in der \texttt{application.properties} konfiguriert werden. Die Token-Validierung erfordert circa 60 Codezeilen manueller Implementierung im \texttt{CognitoService}, einschließlich des Ladens öffentlicher Schlüssel, der Signaturprüfung sowie der Validierung von Issuer und Expiration (vgl. Abbildung \ref{fig:cognito_service}).

    \item \textbf{Imperative Autorisierung ($H_1$, $H_4$):} Autorisierungsregeln werden durch \texttt{@PreAuthorize}-Annotationen auf jeder Controller-Methode definiert (vgl. Abbildung \ref{fig:order_controller}). Dies führt zu einer Verteilung der Sicherheitslogik über den gesamten Quellcode, im Gegensatz zur zentralen Definition im \ac{CDS}-Modell bei \ac{SAP CAP}.

    \item \textbf{Wartbarkeit ($H_2$):} Das Hinzufügen einer neuen Rolle erfordert Änderungen an mehreren \texttt{@PreAuthorize}-Annotationen über verschiedene Controller und Services hinweg. Im Gegensatz zu \ac{SAP CAP} sind die Autorisierungsregeln nicht zentral definiert.

    \item \textbf{Framework-Kopplung ($H_4$):} Die Implementierung ist eng an Spring Security und das \ac{AWS} \ac{SDK} gekoppelt. Mehrere Dateien enthalten \texttt{org.springframework.security} und \texttt{software.amazon.awssdk} Imports, was die Austauschbarkeit des \ac{IdP} erschwert.
\end{itemize}

\section{SAP CAP Framework Backend und SAP IAS}
\label{sec:sap_backend}

Im zweiten Backend wird \ac{SAP IAS} in Kombination mit dem \ac{SAP CAP} Framework mit Java genutzt. Im Gegensatz zum Spring Boot Backend mit \ac{AWS} Cognito verfolgt \ac{SAP CAP} einen deklarativen Ansatz, bei dem Sicherheitsregeln zentral im Datenmodell definiert werden.

\subsection{Architekturübersicht}

Die Architektur des \ac{SAP CAP} Backends unterscheidet sich vom Spring Boot Ansatz:

\begin{itemize}
    \item \textbf{Deklarative Sicherheit:} Autorisierungsregeln werden im \ac{CDS}-Datenmodell definiert.
    \item \textbf{Automatische Token-Validierung:} Das Framework übernimmt die \ac{JWT}-Validierung durch die \texttt{cds-feature-identity} Bibliothek.
    \item \textbf{ODataV4:} Entitäten werden automatisch als OData-Endpoints bereitgestellt.
    \item \textbf{Cloud-Integration:} Die Service-Binding zu XSUAA/\ac{SAP IAS} erfolgt automatisch auf der \ac{SAP BTP}, mithilfe der \texttt{mta.yaml}-Datei.
\end{itemize}

\subsection{Deklarative Autorisierung im CDS-Modell}

Der zentrale Unterschied zu Spring Security liegt in der Definition der Autorisierungsregeln. Abbildung \ref{fig:cds_restrict} zeigt dabei die \texttt{@restrict}-Annotation im Datenmodell.

\begin{figure}[h]
    \begin{verbatim}
entity Orders : cuid, managed {
    customer    : String not null;
    amount      : Decimal(10,2) not null;
    status      : String default 'OPEN';
}

annotate Orders with @(restrict: [
    { grant: 'READ',  to: 'authenticated-user' },
    { grant: ['READ'], to: 'Manager', where: 'createdBy = $user' },
    { grant: ['*'],   to: 'Admin' }
]);
    \end{verbatim}
    \caption{Deklarative Autorisierung durch @restrict-Annotation im CDS-Modell (data-model.cds)}
    \label{fig:cds_restrict}
\end{figure}

Im Vergleich zur imperativen Lösung in Abbildung \ref{fig:order_controller} werden hier alle Autorisierungsregeln an einer zentralen Stelle definiert. Die Regeln umfassen u. a. dabei:
\begin{itemize}
    \item \texttt{authenticated-user}: Jeder authentifizierte Benutzer kann Orders lesen
    \item \texttt{Manager}: Kann eigene Orders lesen
    \item \texttt{Admin}: Vollzugriff auf alle Operationen
\end{itemize}

Die Definition ermöglicht es somit, neue Rollen hinzuzufügen, ohne Java-Code zu ändern oder hinzuzufügen. Diese Regeln dienen zur Absicherung der API-Endpunkte, Service-Endpunkte und Datenbankoperationen.

\subsection{Service-Definition}

Die Service-Definition in \texttt{service-definition.cds} stellt die Entitäten als OData V4 Endpoints bereit.

\begin{figure}[h]
    \begin{verbatim}
using { sap.cap.orders } from '../db/data-model';

service OrderService {
    entity Orders as projection on orders.Orders;

    action completeOrder(orderId: UUID) returns {
        success: Boolean;
        message: String;
    };
}
    \end{verbatim}
    \caption{Service-Definition mit automatischer Autorisierungsdurchsetzung (service-definition.cds)}
    \label{fig:cds_service}
\end{figure}

Die \texttt{@restrict}-Annotationen aus dem Datenmodell werden automatisch auf die Service-Endpoints (Odata) angewendet. Daher ist keine zusätzliche Konfiguration notwendig, da das Framework die Autorisierung durchsetzt.

\subsection{Konfiguration in application.yaml}

Im Gegensatz zur Spring Security Konfiguration für \ac{AWS} Cognito, ist die Konfiguration für \ac{SAP CAP} ebenfalls minimal. Abbildung \ref{fig:cap_config} zeigt die relevanten Einstellungen.

\begin{figure}[h]
    \begin{verbatim}
cds:
  security:
    xsuaa:
      enabled: true
    identity:
      enabled: true
    \end{verbatim}
    \caption{Security-Konfiguration in application.yaml}
    \label{fig:cap_config}
\end{figure}

In der Cloud-Umgebung auf der \ac{SAP BTP} ist zusätzlich die Aktivierung von \ac{XSUAA} erforderlich. Die Verbindung zu \ac{SAP IAS} erfolgt automatisch über Service-Binding. Es ist keine manuelle \ac{JWKS}-\ac{URI} Konfiguration oder JWT-Filter-Implementierung notwendig, da das Framework die gesamte Token-Validierung übernimmt.

\subsection{Event Handler für Geschäftslogik}

Für die Geschäftslogik wird das Setzen des \texttt{createdBy}-Feldes  Event Handler verwendet. Abbildung \ref{fig:cap_handler} zeigt einen Ausschnitt des \texttt{OrderServiceHandler}.

\begin{figure}[h]
    \begin{verbatim}
@Component
@ServiceName(OrderService_.CDS_NAME)
public class OrderServiceHandler implements EventHandler {

    @Before(event = "CREATE", entity = Orders_.CDS_NAME)
    public void beforeCreateOrder(CdsCreateEventContext context) {
        ...
        // Autorisierung erfolgt AUTOMATISCH durch CDS @restrict
        context.getCqn().entries().forEach(entry -> {
            entry.put("createdBy", username);
        });
    }
}
    \end{verbatim}
    \caption{Event-Handler für die Geschäftslogik (\texttt{OrderServiceHandler.java})}
    \label{fig:cap_handler}
\end{figure}

Der Handler enthält keine Autorisierungslogik, diese wird vollständig durch die \texttt{@restrict}-Annotationen im \ac{CDS}-Modell durchgesetzt. Der Handler fokussiert sich ausschließlich auf Geschäftslogik. Im Beispiel wird legiglich, der Nutzername aus den Kontext genommen, um zusätzlich, wenn eine \texttt{Order} erstellt wird, den Ersteller zu dokumentieren.

\subsection{SAP BTP Deployment und Service-Konfiguration}

Für das Deployment auf der \ac{SAP BTP} wird eine Multi-Target Application (MTA) konfiguriert. Die wesentlichen Komponenten sind:

\begin{itemize}
    \item \textbf{XSUAA Service:} Stellt die OAuth2/OIDC-Integration bereit und verbindet sich mit \ac{SAP IAS}
    \item \textbf{App Router:} Handhabt die Authentifizierung und leitet Anfragen an das Backend weiter
    \item \textbf{Backend Service:} Das zuvor implementierte \ac{SAP CAP} Backend
\end{itemize}

Die Rollen werden in der \texttt{xs-security.json} definiert. Abbildung \ref{fig:xs_security} zeigt einen Ausschnitt.

\begin{figure}[h]
    \begin{verbatim}
{
...
  "scopes": [
    { "name": "$XSAPPNAME.Admin", "description": "Full access" },
    { "name": "$XSAPPNAME.Manager", "description": "Read/Write" },
    { "name": "$XSAPPNAME.User", "description": "Read-only" }
  ],
}
    \end{verbatim}
    \caption{Ausschnitt der Rollen-Definition in \texttt{xs-security.json}}
    \label{fig:xs_security}
\end{figure}

\begin{figure}[h]
    \centering
    % TODO: Screenshot der SAP BTP Role Collections Konfiguration
    \fbox{\parbox{0.8\textwidth}{\centering\vspace{2cm}Screenshot: SAP BTP Role Collections\vspace{2cm}}}
    \caption{Role Collections Konfiguration in der SAP BTP (Screenshot folgt)}
    \label{fig:btp_role_collections}
\end{figure}

In Abbildung \ref{fig:btp_role_collections} wird die Zuweisung von Rollen zu Benutzern in der \ac{SAP BTP} dargestellt. Die Konfiguration erfolgt über:
\begin{enumerate}
    \item Erstellung der Role Collections (Admin, Manager, User, Observer)
    \item Zuweisung der Scopes aus \texttt{xs-security.json}
    \item Zuweisung der Benutzer zu den Role Collections
\end{enumerate}

Im Gegensatz zu \ac{AWS} Cognito, bei dem ein User Pool manuell erstellt werden muss (vgl. Abbildung \ref{fig:cognito_user_pool}), erfolgt die \ac{IdP}-Konfiguration bei \ac{SAP IAS} automatisch durch das \ac{XSUAA} Service-Binding in der \texttt{mta.yaml}.

\subsection{Zusammenfassung der Implementierung}

Die Implementierung des \ac{SAP CAP} Backends mit \ac{SAP IAS} zeigt den deklarativen Sicherheitsansatz des Frameworks. Dabei sind folgende Komponente implementiert:

\begin{itemize}
    \item \textbf{CDS-Datenmodell (\texttt{data-model.cds}):} Die Entitätsdefinition für \texttt{Orders} sowie die zentrale Autorisierungslogik durch \texttt{@restrict}-Annotationen. Alle Zugriffsregeln für die Rollen \texttt{authenticated-user}, \texttt{Manager} und \texttt{Admin} sind an dieser zentralen Stelle definiert (vgl. Abbildung \ref{fig:cds_restrict}).

    \item \textbf{Service-Definition (\texttt{service-definition.cds}):} Die Bereitstellung der \texttt{Orders}-Entität als OData V4 Endpoint. Das Framework wendet die \texttt{@restrict}-Annotationen automatisch auf alle generierten Endpunkte an (vgl. Abbildung \ref{fig:cds_service}).

    \item \textbf{Event Handler (\texttt{OrderServiceHandler.java}):} Die Implementierung der Geschäftslogik, die das \texttt{createdBy}-Feld bei der Erstellung einer Order setzt. Der Handler enthält keine Autorisierungslogik, diese wird vollständig durch das \ac{CDS}-Modell durchgesetzt (vgl. Abbildung \ref{fig:cap_handler}).

    \item \textbf{Security-Konfiguration (\texttt{application.yaml}):} Die Aktivierung von \ac{XSUAA} und \ac{SAP IAS} durch zwei Konfigurationszeilen. Keine manuelle \ac{JWKS}-\ac{URI} Konfiguration erforderlich (vgl. Abbildung \ref{fig:cap_config}).

    \item \textbf{MTA-Deployment (\texttt{mta.yaml}):} Die Definition der Cloud-Deployment-Konfiguration mit automatischem Service-Binding zu \ac{XSUAA} und \ac{SAP IAS}.

    \item \textbf{Rollen-Definition (\texttt{xs-security.json}):} Die Definition der Scopes und Role Templates für die \ac{SAP BTP} Role Collections (vgl. Abbildung \ref{fig:xs_security}).
\end{itemize}

\paragraph{Relevanz für die Hypothesen}

\begin{itemize}
    \item \textbf{Automatische Token-Validierung ($H_{3a}$, $H_5$):} Die Token-Validierung erfolgt vollständig durch das Framework. Die \texttt{cds-feature-identity} Bibliothek übernimmt das Laden der \ac{JWKS}-Schlüssel, die Signaturprüfung sowie die Validierung von Issuer und Expiration automatisch. Im Gegensatz zum Spring Boot Backend (vgl. Abbildung \ref{fig:cognito_service}) ist keine manuelle Implementierung erforderlich.

    \item \textbf{Deklarative Autorisierung ($H_1$, $H_4$):} Alle Autorisierungsregeln sind zentral im \ac{CDS}-Modell definiert. Im Vergleich zum imperativen Ansatz mit \texttt{@PreAuthorize}-Annotationen (vgl. Abbildung \ref{fig:order_controller}) befinden sich die Sicherheitsregeln nicht verteilt über Controller und Services, sondern an einer zentralen Stelle.

    \item \textbf{Wartbarkeit ($H_2$):} Das Hinzufügen einer neuen Rolle erfordert lediglich eine Erweiterung der \texttt{@restrict}-Annotation in der \texttt{data-model.cds}. Keine Änderungen in Java-Controllern, Services oder Event Handlern sind notwendig.
\end{itemize}
